% ============================================================
%  CAPÍTULO 9 – Planificación, Diagnóstico y Contexto Laboral
% ============================================================
\chapter{Planificación, Diagnóstico y Contexto Laboral}

\section{Planificación del proyecto -- Diagrama de Gantt}

El desarrollo de FitLabs se planificó en seis fases. A continuación se muestra el diagrama de Gantt estimado con las duraciones reales de cada fase:

\begin{figure}[H]
\centering
\resizebox{0.95\textwidth}{!}{\includegraphics{gantt_chart.png}}
\caption{Diagrama de Gantt del proyecto FitLabs (octubre 2025 -- mayo 2026).}
\label{fig:gantt}
\end{figure}

\section{Análisis DAFO}

El análisis DAFO permite evaluar los factores internos y externos que afectan al desarrollo y viabilidad de FitLabs:

\begin{figure}[H]
  \centering
  \begin{tikzpicture}[every node/.style={draw, rounded corners, inner sep=8pt, text width=0.42\textwidth}, node distance=6mm]
    \node (F) [fill=green!10] {\textbf{Fortalezas}\\- Interfaz intuitiva y diseño centrado en el usuario\\- Desarrollo con Flutter: rendimiento y UI consistente\\- Integración con servicios modernos (Supabase, APIs)};
    \node (D) [right=of F, fill=red!5] {\textbf{Debilidades}\\- Equipo y recursos limitados\\- Dependencia de servicios externos\\- Curva de aprendizaje en ciertas integraciones};
    \node (O) [below=of F, fill=blue!5] {\textbf{Oportunidades}\\- Mercado creciente de apps de fitness y salud\\- Integración con wearables y datos biométricos\\- Colaboraciones con gimnasios y profesionales};
    \node (A) [right=of O, fill=yellow!8] {\textbf{Amenazas}\\- Competencia consolidada (apps comerciales)\\- Cambios en políticas de APIs o precios\\- Riesgos de seguridad y privacidad de datos};
  \end{tikzpicture}
  \caption{Análisis DAFO (SWOT) resumido para el proyecto FitLabs}
  \label{fig:dafo_fitlabs}
\end{figure}
\FloatBarrier

\begin{figure}[H]
  \centering
  \begin{tikzpicture}
    \begin{polaraxis}[
      width=0.75\textwidth,
      height=7cm,
      ymin=0,ymax=10,
      xtick={0,90,180,270},
      xticklabels={Fortalezas,Oportunidades,Debilidades,Amenazas},
      ytick={0,2,4,6,8,10},
      grid=both,
      major grid style={gray!30},
      minor grid style={gray!10},
    ]
    \addplot+[mark=none,fill=blue!30,fill opacity=0.6] coordinates {(0,8) (90,7) (180,4) (270,5) (0,8)};
    \end{polaraxis}
  \end{tikzpicture}
  \caption{Gráfico DAFO: puntuaciones de ejemplo (0--10).}
  \label{fig:dafo_grafico}
\end{figure}
\FloatBarrier
